\par
\begingroup
\renewcommand{\arraystretch}{1.45}
\SourceLeaf{282}
Multiplicentur hae aequationes resp. per
\[
\begin{gathered}
\frac{1}{(a-b)(a-c)(a-d)\mathinner{.\,.\,.\,.}(a-t)}\\[4pt]
\frac{1}{(b-a)(b-c)(b-d)\mathinner{.\,.\,.\,.}(b-t)}\\[4pt]
\frac{1}{(c-a)(c-b)(c-d)\mathinner{.\,.\,.\,.}(c-t)}\\[4pt]
\frac{1}{(d-a)(d-b)(d-c)\mathinner{.\,.\,.\,.}(d-t)}\\
\text{etc.}\\
\frac{1}{(t-a)(t-b)(t-c)\mathinner{.\,.\,.\,.}}
\end{gathered}
\]
prodeatque inde per productorum additionem
\[
\begin{aligned}
&\frac{A}{(a-b)(a-c)(a-d)\mathinner{.\,.\,.}(a-t)}\\
&{}+\frac{B}{(b-a)(b-c)(b-d)\mathinner{.\,.\,.}(b-t)}\\
&{}+\frac{C}{(c-a)(c-b)(c-d)\mathinner{.\,.\,.}(c-t)}\\
&{}+\frac{D}{(d-a)(d-b)(d-c)\mathinner{.\,.\,.}(d-t)}\\
&{}+\text{etc.}\\
&{}+\frac{T}{(t-a)(t-b)(t-c)(t-d)\mathinner{.\,.\,.}}=W
\end{aligned}
\]
Tunc ex art. 1, ubi $m$ idem denotabat, quod hic nobis est $m+1$, facile concludetur, fieri $W=0$; quamobrem multiplicando per $(t-a)(t-b)(t-c)\mathinner{.\,.\,.}$ prodit
\[
\begin{aligned}
T={}&\frac{(t-b)(t-c)(t-d)\mathinner{.\,.\,.}}{(a-b)(a-c)(a-d)\mathinner{.\,.\,.}}A\\
&{}+\frac{(t-a)(t-c)(t-d)\mathinner{.\,.\,.}}{(b-a)(b-c)(b-d)\mathinner{.\,.\,.}}B\\
&{}+\frac{(t-a)(t-b)(t-d)\mathinner{.\,.\,.}}{(c-a)(c-b)(c-d)\mathinner{.\,.\,.}}C\\
&{}+\frac{(t-a)(t-b)(t-c)\mathinner{.\,.\,.}}{(d-a)(d-b)(d-c)\mathinner{.\,.\,.}}D\\
&{}+\text{etc.}
\end{aligned}
\]
\InterpolationArticle{4}
Formula in art. praec. inventa, ita comparata est, ut sponte sine omni calculo pateat, si pro $t$ quantitatum $a,b,c,d\mathinner{.\,.}$ aliqua in illa substituatur, valorem
respondentem $A,B,C,D\mathinner{.\,.}$ inde prodire. Neque hoc solo respectu sese commendat: certo enim ad usum practicum longe commodissima est, saltem quoties uni\SourceLeaf{283}%
cus tantum valor ipsius $X$ e valoribus datis computandus est. Quando plures
sunt eruendi, transformatio sequens formulae nostrae nonnunquam praeferri
poterit
\[
\begin{aligned}
T={}&A+(t-a)\left(\frac{A}{a-b}+\frac{B}{b-a}\right)\\
&{}+(t-a)(t-b)\left(\frac{A}{(a-b)(a-c)}+\frac{B}{(b-a)(b-c)}+\frac{C}{(c-a)(c-b)}\right)\\
&{}+(t-a)(t-b)(t-c)\left(\frac{A}{(a-b)(a-c)(a-d)}+\frac{B}{(b-a)(b-c)(b-d)}\right.\\
&\qquad\left.{}+\frac{C}{(c-a)(c-b)(c-d)}+\frac{D}{(d-a)(d-b)(d-c)}\right)\\
&{}+\text{etc.}
\end{aligned}
\]
(quae ex formula art. 3 facillime derivatur, si in hac multitudo quantitatum
$a,b,c,d\mathinner{.\,.\,.}$ ab una $a$ ad duas $a,b$, inde ad tres $a,b,c$ etc. successive increscere
concipitur et quisque valor ipsius $T$ hoc modo oriens a sequenti subtrahitur).
Ponantur coëfficientes
\[
\begin{gathered}\frac{A}{a-b}+\frac{B}{(b-a)},\quad
\frac{A}{(a-b)(a-c)}+\frac{B}{(b-a)(b-c)}+\frac{C}{(c-a)(c-b)}\text{ etc.}\\
=A',\ A''\text{ etc.}\end{gathered}
\]
porro designetur per $B',B''$ etc. id, quod fit ex $A',A''$ etc. si pro $a,b,c\mathinner{.\,.\,.}$
resp. scribitur $b,c,d\mathinner{.\,.\,.}$ atque pro $A,B,C\mathinner{.\,.\,.}$ resp. $B,C,D\mathinner{.\,.\,.}$; similiter designetur per $C',C''$ etc. id quod fit ex $A',A''$ etc., si pro $a,b,c\mathinner{.\,.\,.}$ resp. substituitur $c,d,e\mathinner{.\,.\,.}$ et pro $A,B,C\mathinner{.\,.\,.}$ resp. $C,D,E\mathinner{.\,.\,.}$ et sic porro. Tunc habebimus
\[
\begin{array}{lll}
A'=\dfrac{A-B}{a-b},&B'=\dfrac{B-C}{b-c},&C'=\dfrac{C-D}{c-d}\ \mathinner{.\,.\,.}\\[5pt]
A''=\dfrac{A'-B'}{a-c},&B''=\dfrac{B'-C'}{b-d},&C''=\dfrac{C'-D'}{c-e}\ \mathinner{.\,.\,.}\\[5pt]
A'''=\dfrac{A''-B''}{a-d},&B'''=\dfrac{B''-C''}{b-e},&C'''=\dfrac{C''-D''}{c-f}\ \mathinner{.\,.\,.}\\
\text{etc.}&&
\end{array}
\]
atque
\[\begin{aligned}
T&=A+A'(t-a)+A''(t-a)(t-b)\\&+A'''(t-a)(t-b)(t-c)+\mathinner{.\,.}
\end{aligned}\]
Multitudo quantitatum $A',B',C'$ etc. hic computandarum erit $m-1$, multitudo
quantitatum $A'',B'',C''$ etc. erit $m-2$ et sic porro.
\InterpolationArticle{5}
Sequens quoque mutatio formulae art. 3 non sine fructu in usum vocari
poterit. Sit $\mathfrak{X}$ functio integra arbitraria indeterminatae $x$, neque tamen altioris
\SourceLeaf{284}
quam $m-1^{\mathrm{ti}}$ ordinis; atque $\mathfrak{A},\mathfrak{B},\mathfrak{C}.\mathfrak{D}\mathinner{.\,.\,.\,.}\mathfrak{T}$ eius valores pro $x=a,b,c,d\mathinner{.\,.\,.}t$
resp. Tunc per art. 1 erit
\[
\begin{aligned}
0={}&\frac{\mathfrak{A}}{(a-b)(a-c)(a-d)\mathinner{.\,.\,.}(a-t)}\\
&{}+\frac{\mathfrak{B}}{(b-a)(b-c)(b-d)\mathinner{.\,.\,.}(b-t)}\\
&{}+\frac{\mathfrak{C}}{(c-a)(c-b)(c-d)\mathinner{.\,.\,.}(c-t)}\\
&{}+\frac{\mathfrak{D}}{(d-a)(d-b)(d-c)\mathinner{.\,.\,.}(d-t)}\\
&{}+\text{etc.}\\
&{}+\frac{\mathfrak{T}}{(t-a)(t-b)(t-c)(t-d)\mathinner{.\,.\,.}}
\end{aligned}
\]
Hinc facile deducitur
\[
\begin{aligned}
T={}&\mathfrak{T}\\
&{}+\frac{(t-b)(t-c)(t-d)\mathinner{.\,.\,.}}{(a-b)(a-c)(a-d)\mathinner{.\,.\,.}}(A-\mathfrak{A})\\
&{}+\frac{(t-a)(t-c)(t-d)\mathinner{.\,.\,.}}{(b-a)(b-c)(b-d)\mathinner{.\,.\,.}}(B-\mathfrak{B})\\
&{}+\frac{(t-a)(t-b)(t-d)\mathinner{.\,.\,.}}{(c-a)(c-b)(c-d)\mathinner{.\,.\,.}}(C-\mathfrak{C})\\
&{}+\frac{(t-a)(t-b)(t-c)\mathinner{.\,.\,.}}{(d-a)(d-b)(d-c)\mathinner{.\,.\,.}}(D-\mathfrak{D})\\
&{}+\text{etc.}
\end{aligned}
\]
Quodsi functio $\mathfrak{X}$ ita eligitur, ut valores functionis $X$ hi $A,B,C,D\mathinner{.\,.\,.}$ prope
per illam repraesententur, hoc lucramur, ut coëfficientes
$\frac{(t-b)(t-c)(t-d)\mathinner{.\,.\,.}}{(a-b)(a-c)(a-d)\mathinner{.\,.\,.}}$ etc.
per quantitates parvas multiplicandi sint. Potest etiam pro $\mathfrak{X}$ quantitas constans
assumi, e. g. una ex his $A,B,C,D\mathinner{.\,.\,.}$; in quo casu pars una e valore ipsius
$T$ excidet. Aut ita determinari potest $\mathfrak{X}$, ut duae pluresve harum quantitatum
per $\mathfrak{X}$ repraesententur, in quo casu totidem partes ex $T$ excident.
\InterpolationArticle{6}
Quae praecedunt, suppositioni innituntur, functionem $X$ ultra potestatem
$x^{m-1}$ non egredi, sive differentiam $m^{\mathrm{tam}}$ una cum superioribus evanescere, in
quo casu methodus interpolationis rigorose vera est. Si vero illa suppositio locum
non habet, interpolatio eo tantummodo tendit, ut loco functionis $X$ functio \textit{simplicissima} eruatur, per quam valoribus propositis $A,B,C,D\mathinner{.\,.\,.}$ satisfiat. Iam
ut errorem, qui a neglectis differentiis superioribus nascitur, diiudicare possimus,
\SourceLeaf{285}
sint termini in \(X\) post potestatem \(x^{m-1}\) hi \(\mu x^m+\nu x^{m+1}+\text{etc.}\) Erit itaque in art. 3
non \(W=0\), sed ut ex art. 1 sequitur, \(W=\mu+\nu(a+b+c+d+\mathinner{.\,.\,.}+t)+\text{etc.}\)
Hinc valori ipsius \(T\) illic tradito adhuc adiici debet
\[
(t-a)(t-b)(t-c)(t-d)\mathinner{.\,.\,.}\times\{\mu+\nu(a+b+c+d+\mathinner{.\,.\,.}+t)+\text{etc.}\}
\]

\InterpolationArticle{7}
Casus in praxi maxime frequens est, ubi \(a,b,c,d\mathinner{.\,.\,.}\) progressionem arithmeticam constituunt. Ponendo intervallum \(=1\), ita ut sit \(b=a+1,\ c=a+2\)
etc., formula art. 4 fit
\[
\begin{aligned}
T={}&A\\
&+(t-1)(B-A)\\
&+(t-1)(t-2)\left(\frac{C-2B+A}{2}\right)\\
&+(t-1)(t-2)(t-3)\left(\frac{D-3C+3B-A}{6}\right)\\
&+\text{etc.}
\end{aligned}
\]
sive
\[\begin{aligned}
T&=A+A'(t-1)+A''(t-1)(t-2)\\&+A'''(t-1)(t-2)(t-3)+\mathinner{.\,.}
\end{aligned}\]
ubi \(A',A'',A'''\) etc. computantur per algorithmum sequentem
\[
\begin{array}{lll}
A'=B-A,&B'=C-B,&C'=D-C\quad\text{etc.}\\
2A''=B'-A',&2B''=C'-B',&2C''=D'-C'\quad\text{etc.}\\
3A'''=B''-A'',&3B'''=C''-B'',&3C'''=D''-C''\quad\text{etc.}
\end{array}
\]
etc., quae formula cum vulgata interpolationis formula per differentias omnino
convenit.

\InterpolationArticle{8}
Si pro satis multis valoribus ipsius \(x\) in serie arithmetica progredientibus
\(a,a+1,a+2\mathinner{.\,.\,.}\) valores respondentes functionis \(X\) cogniti sunt, ut seriem \(m\)
valorum successivorum ipsius \(x\) ad libitum eligere liceat ad computum ipsius \(T\):
quaestio oritur, \emph{quosnam} valores ad hunc finem praeferre maxime praestet, siquidem plures quam \(m\) adhibere sive ultra differentiam \(m-1^{\mathrm{tam}}\) egredi nolimus?
\SourceLeaf{286}
Manifesto hoc ita est decidendum, ut error a differentia \(m^{\mathrm{ti}}\) ordinis oriundus fiat
quam minimus; hic vero error fit
\[
=(t-a)(t-a-1)(t-a-2)\mathinner{.\,.\,.\,.}(t-a-m+1)\frac{\Delta}{1.2.3\mathinner{.\,.\,.} m}
\]
si adhibentur termini ad \(x=a,a+1,a+2,\mathinner{.\,.\,.} a+m-1\) pertinentes, vel
\[
=(t-a-1)(t-a-2)(t-a-3)\mathinner{.\,.\,.\,.}(t-a-m)\frac{\Delta}{1.2.3\mathinner{.\,.\,.} m}
\]
si adhibentur termini ad \(x=a+1,a+2,a+3,\mathinner{.\,.\,.} a+m\) pertinentes, vel
\[
=(t-a+1)(t-a)(t-a-1)\mathinner{.\,.\,.}(t-a-m+2)\frac{\Delta}{1.2.3\mathinner{.\,.\,.\,.}m}
\]
si adhibentur termini ad \(x=a-1,a,a+1,\mathinner{.\,.\,.} a+m-2\) pertinentes, designando
per \(\Delta\) differentiam ordinis \(m^{\mathrm{ti}}\), quae siquidem ad differentias superiorum ordinum
non respicitur, hic tamquam constans considerari potest. Quodsi igitur \(m\) termini
ad \(x=a,a+1,a+2,\mathinner{.\,.\,.} a+m-1\) pertinentes maxime idonei sunt, ex illis
tribus expressionibus prima debet esse minima adeoque sine respectu signi \(t-a<\)
vel saltem non \(>t-a-m\), nec non \(t-a+1>\) vel saltem non \(<t-a-m+1\).
Ex conditione priore facile deducitur, \(t-a-m\) esse debere quantitatem negativam inter \(-\frac12m\) et \(-\infty\) sitam, unde \(t-a\) iacebit inter \(+\frac12m\) et \(-\infty\); ex
conditione posteriore autem erit \(t-a+1\) quantitas positiva atque inter \(\frac12m\) et
\(+\infty\) sita; quare \(t-a\) iacebit inter \(\frac12m-1\) et \(\frac12m\). Hinc sequitur, si \(m\) fuerit numerus par, valores ad interpolationem adhibendos ita eligendos esse, ut
prior semissis ab una parte, posterior ab altera termini quaesiti iaceant; si autem
\(m\) fuerit impar, termini ii sunt eligendi, quorum medius quaesito iaceat quam
proximus. Quoties in hoc casu \(t\) est exacte medius inter duos valores consecutivos ipsius \(x\), respectu erroris a neglecta differentia \(m^{\mathrm{ti}}\) ordinis generaliter loquendo nihil interest, sive terminus quaesito praecedens sive sequens pro medio
adhibendorum adoptetur.

\InterpolationArticle{9}
\emph{Exempla.} I. Invenire oporteat \(\log.\sin 24^{0}30'\), si logarithmi sinuum per singulos gradus habentur, ita ut differentiae quartae et superiores negligantur. In
hoc itaque casu secundum praecepta art. praec. adhibebimus logarithmos sinuum
arcuum \(23^{0},24^{0},25^{0},26^{0}\), quibus per \(A,B,C,D\) designatis provenit per formulam art. 3 logarithmus quaesitus
\SourceLeaf{287}
\[
=-\frac{1}{16}A+\frac{9}{16}B+\frac{9}{16}C-\frac{1}{16}D
\]
cuius valor numericus ex
\[
\begin{aligned}
A&=9{,}5918780\\
B&=9{,}6093133\\
C&=9{,}6259483\\
D&=9{,}6418420
\end{aligned}
\]
eruitur \(9{,}6177271.5\); tabulae dant \(9{,}6177270\).

Quodsi hic loco \(\log.\sin 23^{0}\), logarithmum \(\sin 27^{0}=E=9{,}6570468\) adhibuissemus, formula
\[
+\frac{5}{16}B+\frac{15}{16}C-\frac{5}{16}D+\frac{1}{16}E
\]
dedisset \(9{,}6177267\,.\,375\).

Formula art. praec. pro hoc casu dat errorem ex neglecta differentia quarta
\(\Delta\) oriundum \(=\frac{3}{128}\Delta\), si \(A,B,C,D\), contra \(=-\frac{5}{128}\Delta\), si \(B,C,D\)\raisebox{-0.3ex}{:}\,\(E\) adhibentur; qui error valori calculato adiiciendus est. In exemplo nostro habetur
\(\Delta=-0{,}0000066\); correctio hinc calculata utrumque computum cum tabulis
conciliat.

II. Sit tempus novilunii 16 Junii 1806 e longitudinibus solis et lunae pro
singulis diebus computandum, ita ut differentiarum usque ad quartam incl. ratio
habeatur. Habentur in ephemeridibus Parisiensibus pro meridie vero loci sequentes

\begin{center}
\begin{tabular}{r@{\qquad}rrr@{\qquad}rrr}
&\multicolumn{3}{c}{longit. solis}&\multicolumn{3}{c}{longit. lunae}\\
1806 Junii 14&\(82^{0}\)&\(39'\)&\(12''\)&\(53^{0}\)&\(19'\)&\(15''\)\\
15&83&36&30&67&29&43\\
16&84&33&48&81&59&52\\
17&85&31&7&96&44&9\\
18&86&28&24&111&35&30
\end{tabular}
\end{center}

Sunt hic valores dati ipsius \(x\) differentiae inter longitudinem solis et lunae, valores respondentes ipsius \(X\) tempora, quae inde a meridie 16 Junii numerabuntur; quaeriturque valor ipsius \(X\) valori \(x=0\) respondens. Fit igitur
\SourceLeaf{288}
\[
\begin{array}{r@{\;}c@{\;}r@{\;}r@{\;}r@{\quad}c@{\;}r@{\qquad}r@{\;}c@{\;}r}
a&=&-29^{0}&19'&57''&=&-105596''&A&=&-2\\
b&=&-16&6&47&=&-58007&B&=&-1\\
c&=&-2&33&56&=&-9236&C&=&0\\
d&=&+11&13&2&=&+40382&D&=&+1\\
e&=&+25&7&6&=&+90426&E&=&+2\\
t&=&0&&&&&&&
\end{array}
\]
Hinc computatur
\[
\begin{aligned}
\frac{(t-b)(t-c)(t-d)(t-e)}{(a-b)(a-c)(a-d)(a-e)}&=+0{,}0149084\\[5pt]
\frac{(t-a)(t-c)(t-d)(t-e)}{(b-a)(b-c)(b-d)(b-e)}&=-0{,}1050660\\[5pt]
\frac{(t-a)(t-b)(t-d)(t-e)}{(c-a)(c-b)(c-d)(c-e)}&=+0{,}9624567\\[5pt]
\frac{(t-a)(t-b)(t-c)(t-e)}{(d-a)(d-b)(d-c)(d-e)}&=+0{,}1434435\\[5pt]
\frac{(t-a)(t-b)(t-c)(t-d)}{(e-a)(e-b)(e-c)(e-d)}&=-0{,}0157429
\end{aligned}
\]
His coëfficientibus\HistoricalNote{23} per \(-2,-1,0,1,2\) resp. multiplicatis, productorum summa
fit \(=+0{,}1872069\); quare novilunium erit 16 Junii, \(4^{\mathrm h}29'34''7\) temp. ver.
Paris. sive \(4^{\mathrm h}29'41''4\) temp. med.; Connaissance des tems habet \(4^{\mathrm h}27'42''\).

\InterpolationArticle{10}
Sit \(X\) functio arcus indeterminati \(x\) huius formae
\[
\begin{aligned}
&\alpha+\alpha'\cos x+\alpha''\cos 2x+\alpha'''\cos 3x+\text{etc.}\\
&\quad+\beta'\sin x+\beta''\sin 2x+\beta'''\sin 3x+\text{etc.}
\end{aligned}
\]
quae non excurrat in infinitum, sed cum \(\cos mx\) et \(\sin mx\) abrumpatur, ita
ut multitudo coëfficientium (incognitorum) sit \(2m+1\). Pro totidem valoribus diversis ipsius \(x\), puta \(a,b,c,d\) etc. dati sint valores respondentes functionis \(X\),
puta \(A,B,C,D\mathinner{.\,.\,.}\) (Ceterum valores ipsius \(x\), quorum differentia est peripheria integra sive eius multiplum, manifesto hic pro diversis haberi nequeunt). Ex
his datis quaeritur formula pro valore \(T\), quem functio \(X\) pro quocunque alio
valore ipsius \(x\), puta \(t\) nanciscitur. Habentur itaque \(2m+2\) aequationes
\SourceLeaf{289}
\[
\begin{aligned}
A&=\alpha+\alpha'\cos a+\alpha''\cos 2a+\alpha'''\cos 3a+\text{etc.}\\
 &\quad+\beta'\sin a+\beta''\sin 2a+\beta'''\sin 3a+\text{etc.}\\
B&=\alpha+\alpha'\cos b+\alpha''\cos 2b+\alpha'''\cos 3b+\text{etc.}\\
 &\quad+\beta'\sin b+\beta''\sin 2b+\beta'''\sin 3b+\text{etc.}\\
C&=\alpha+\alpha'\cos c+\alpha''\cos 2c+\alpha'''\cos 3c+\text{etc.}\\
 &\quad+\beta'\sin c+\beta''\sin 2c+\beta'''\sin 3c+\text{etc.}\\
D&=\alpha+\alpha'\cos d+\alpha''\cos 2d+\alpha'''\cos 3d+\text{etc.}\\
 &\quad+\beta'\sin d+\beta''\sin 2d+\beta'''\sin 3d+\text{etc.}\\
 &\text{etc.}\\
T&=\alpha+\alpha'\cos t+\alpha''\cos 2t+\alpha'''\cos 3t+\text{etc.}\\
 &\quad+\beta'\sin t+\beta''\sin 2t+\beta'''\sin 3t+\text{etc.}
\end{aligned}
\]
Multiplicando has aequationes resp. per
\[
\begin{gathered}
\frac{1}{\sin\tfrac12(a-b)\sin\tfrac12(a-c)\sin\tfrac12(a-d)\mathinner{.\,.\,.\,.}\sin\tfrac12(a-t)}\\[6pt]
\frac{1}{\sin\tfrac12(b-a)\sin\tfrac12(b-c)\sin\tfrac12(b-d)\mathinner{.\,.\,.\,.}\sin\tfrac12(b-t)}\\[6pt]
\frac{1}{\sin\tfrac12(c-a)\sin\tfrac12(c-b)\sin\tfrac12(c-d)\mathinner{.\,.\,.\,.}\sin\tfrac12(c-t)}\\[6pt]
\frac{1}{\sin\tfrac12(d-a)\sin\tfrac12(d-b)\sin\tfrac12(d-c)\mathinner{.\,.\,.\,.}\sin\tfrac12(d-t)}\\
\text{etc.}\\[6pt]
\frac{1}{\sin\tfrac12(t-a)\sin\tfrac12(t-b)\sin\tfrac12(t-c)\sin\tfrac12(t-d)\mathinner{.\,.\,.}}
\end{gathered}
\]
addendoque producta prodeat
\[
\begin{aligned}
&\frac{A}{\sin\tfrac12(a-b)\sin\tfrac12(a-c)\sin\tfrac12(a-d)\mathinner{.\,.\,.\,.}\sin\tfrac12(a-t)}\\[6pt]
+{}&\frac{B}{\sin\tfrac12(b-a)\sin\tfrac12(b-c)\sin\tfrac12(b-d)\mathinner{.\,.\,.\,.}\sin\tfrac12(b-t)}\\[6pt]
+{}&\frac{C}{\sin\tfrac12(c-a)\sin\tfrac12(c-b)\sin\tfrac12(c-d)\mathinner{.\,.\,.\,.}\sin\tfrac12(c-t)}\\[6pt]
+{}&\frac{D}{\sin\tfrac12(d-a)\sin\tfrac12(d-b)\sin\tfrac12(d-c)\mathinner{.\,.\,.\,.}\sin\tfrac12(d-t)}\\
+{}&\text{etc.}\\[6pt]
+{}&\frac{T}{\sin\tfrac12(t-a)\sin\tfrac12(t-b)\sin\tfrac12(t-c)\sin\tfrac12(t-d)\mathinner{.\,.\,.\,.}}=W
\end{aligned}
\]
\SourceLeaf{290}
Tunc ex art. 2, ubi $m$ idem erat, quod hic est $2m+2$, adeoque casus secundus locum habet, fit $W=0$; quamobrem multiplicando per
\[
\sin\tfrac12(t-a)\sin\tfrac12(t-b)\sin\tfrac12(t-c)\sin\tfrac12(t-d)\mathinner{.\,.\,.}
\]
prodit
\[
\begin{aligned}
T={}&\frac{\sin\tfrac12(t-b)\sin\tfrac12(t-c)\sin\tfrac12(t-d)\mathinner{.\,.\,.\,.}}
{\sin\tfrac12(a-b)\sin\tfrac12(a-c)\sin\tfrac12(a-d)\mathinner{.\,.\,.\,.}}A\\[5pt]
&+\frac{\sin\tfrac12(t-a)\sin\tfrac12(t-c)\sin\tfrac12(t-d)\mathinner{.\,.\,.\,.}}
{\sin\tfrac12(b-a)\sin\tfrac12(b-c)\sin\tfrac12(b-d)\mathinner{.\,.\,.\,.}}B\\[5pt]
&+\frac{\sin\tfrac12(t-a)\sin\tfrac12(t-b)\sin\tfrac12(t-d)\mathinner{.\,.\,.\,.}}
{\sin\tfrac12(c-a)\sin\tfrac12(c-b)\sin\tfrac12(c-d)\mathinner{.\,.\,.\,.}}C\\[5pt]
&+\frac{\sin\tfrac12(t-a)\sin\tfrac12(t-b)\sin\tfrac12(t-c)\mathinner{.\,.\,.\,.}}
{\sin\tfrac12(d-a)\sin\tfrac12(d-b)\sin\tfrac12(d-c)\mathinner{.\,.\,.\,.}}D\\[5pt]
&+\text{etc.}
\end{aligned}
\]
Quum haec formula indefinite pro valore quocunque ipsius $t$ locum habeat, manifestum est, si producta sinuum in numeratoribus in cosinus sinusque arcuum
multiplicium evolvantur, id quod provenit cum
\[
\begin{aligned}
&\alpha+\alpha'\cos t+\alpha''\cos2t+\alpha'''\cos3t+\text{etc.}\\
&\quad+\beta'\sin t+\beta''\sin2t+\beta'''\sin3t+\text{etc.}
\end{aligned}
\]
\emph{identicum} esse debere, unde coëfficientes $\alpha,\alpha',\beta',\alpha'',\beta''$ etc. innotescent. Ceterum formula pro $T$, ut hic exhibita est, ita est comparata, ut sponte et sine calculo pateat, substitutis pro $t$ resp. $a,b,c,d$ etc. valoribus propositis $A,B,C,D$
etc. probe satisfieri.

\InterpolationArticle{11}

Si multitudo valorum datorum $A,B,C,D$ etc. unitate minor esset, quam
supposuimus, puta $=2m$, hi non sufficient ad determinationem $2m+1$ coëfficientium incognitorum, nisi inter eos relatio aliqua cognita fuerit. In hoc itaque casu expressio generalis pro $T$ aliter exhiberi nequit, nisi ita, ut constantem incognitam implicet. Ratiocinia art. praec. hic adhiberi non possunt, quoniam summatio $W=0$, ex art. 2 petita, suppositioni innixa est, multitudinem
quantitatum $A,B\mathinner{.\,.\,.}T$ esse parem. Formula quidem in art. praec. pro $T$ inventa valoribus $A,B$ etc. generaliter satisfacit, sive horum multitudo par sit
sive impar: sed levi attentione perspicitur, ex evolutione productorum e sinubus
prodire, in nostro casu, expressionem propositae non similem, sed huius formae
\SourceLeaf{291}
\[
\begin{aligned}
&\mathfrak{A}\cos\tfrac12t+\mathfrak{A}'\cos\tfrac32t+\mathfrak{A}''\cos\tfrac52t+\text{etc.}\\
&\quad+\mathfrak{B}\sin\tfrac12t+\mathfrak{B}'\sin\tfrac32t+\mathfrak{B}''\sin\tfrac52t+\text{etc.}
\end{aligned}
\]
Quamobrem pro hoc casu viam aliam ingredi oportebit.

Multiplicetur $X$ per $\cos(\tfrac12x+k)$, proditque, si termini ultimi in $X$ statuuntur $\alpha^m\cos mx+\beta^m\sin mx$
\[
\begin{aligned}
&\alpha\cos(\tfrac12x+k)+\tfrac12\alpha'\cos(\tfrac32x+k)\\&\qquad+\tfrac12\alpha''\cos(\tfrac52x+k)+\mathinner{.\,.\,.\,.}\\
&\hspace{2.01em}+\tfrac12\alpha^m\cos((m+\tfrac12)x+k)\\
&+\tfrac12\alpha'\cos(\tfrac12x-k)+\tfrac12\alpha''\cos(\tfrac32x-k)\\&\qquad+\tfrac12\alpha'''\cos(\tfrac52x-k)+\mathinner{.\,.\,.}\\
&\hspace{2.01em}+\tfrac12\alpha^m\cos((m-\tfrac12)x-k)\\
&\hspace{2em}+\tfrac12\beta'\sin(\tfrac32x+k)+\tfrac12\beta''\sin(\tfrac52x+k)+\mathinner{.\,.\,.}\\
&\hspace{2.01em}+\tfrac12\beta^m\sin((m+\tfrac12)x+k)\\
&+\tfrac12\beta'\sin(\tfrac12x-k)+\tfrac12\beta''\sin(\tfrac32x-k)\\&\qquad+\tfrac12\beta'''\sin(\tfrac52x-k)+\mathinner{.\,.\,.}\\
&\hspace{2.01em}+\tfrac12\beta^m\sin((m-\tfrac12)x-k)
\end{aligned}
\]
Iam si in hac expressione pro $x$ successive substituitur $a,b,c,d\mathinner{.\,.\,.\,.}t$, quod
inde provenit resp. per
\[
\begin{aligned}
&\frac{1}{\sin\tfrac12(a-b)\sin\tfrac12(a-c)\sin\tfrac12(a-d)\mathinner{.\,.\,.\,.}\sin\tfrac12(a-t)}\\[1ex]
&\frac{1}{\sin\tfrac12(b-a)\sin\tfrac12(b-c)\sin\tfrac12(b-d)\mathinner{.\,.\,.\,.}\sin\tfrac12(b-t)}\\[1ex]
&\frac{1}{\sin\tfrac12(c-a)\sin\tfrac12(c-b)\sin\tfrac12(c-d)\mathinner{.\,.\,.\,.}\sin\tfrac12(c-t)}\\[1ex]
&\frac{1}{\sin\tfrac12(d-a)\sin\tfrac12(d-b)\sin\tfrac12(d-c)\mathinner{.\,.\,.\,.}\sin\tfrac12(d-t)}\\
&\text{etc.}\\[1ex]
&\frac{1}{\sin\tfrac12(t-a)\sin\tfrac12(t-b)\sin\tfrac12(t-c)\sin\tfrac12(t-d)\mathinner{.\,.\,.\,.}}
\end{aligned}
\]
multiplicatur, productorumque aggregatum $=W$ ponitur: ex iis, quae in art. 2,
(in casu priori, ubi $m$ idem erat, quod hic est $2m+1$) tradidimus, facile sequitur, haberi
\[
\begin{aligned}
W={}&\tfrac12(2i)^{2m}\alpha^m\cos(\tfrac12(a+b+c+d+\mathinner{.\,.\,.}+t)+k)\\
&+\tfrac12(2i)^{2m}\beta^m\sin(\tfrac12(a+b+c+d+\mathinner{.\,.\,.}+t)+k)
\end{aligned}
\]
Ex altera vero parte erit
\endgroup
